\part{QoS-targeted control: making the slices meet their SLAs on the go}

\begin{keybox}{Goal of this part}
After the engineering audit, the goal was restated as: \emph{SliceNet and RL should slice the network so that the four
classes each get the bandwidth they need to meet their QoS requirements, adapting on the fly, predictively or
reactively, but effectively.} This part records what stood in the way, the six modules built to close the gap
(A–F, Figure~\ref{fig:qosloop}), and what the measurements say. Every module was regression-tested on the emulated
network before the next one started. Development numbers marked \emph{(dev, $n{=}2$)} come from two-seed checks
used to steer the design; the verdicts come from the contended-regime evaluation (\S\ref{sec:qoseval}).
\end{keybox}

\FigQosLoop

\section{Diagnosis: why the audited system could not meet the goal}

\begin{qa}{701}{In the audited system, where did the classes actually miss their SLAs?}
Only in the emergency contexts, and only on loss. In the main study, healthcare met its SLA in every window in normal,
school and harvest. In disaster and medical emergency it failed in 22\,\% and 33\,\% of windows. Delay and jitter never
failed; loss above 1\,\% did, all of it from telemedicine video (3.5–6\,\% loss). ECG and voice lost nothing, because
fq\_codel's sparse-flow priority already protects small flows inside a slice.
\ev{\Lthree}{\code{experiments/qos_compare.py} on the main-study runs; per-application loss from the measured datasets.}
\end{qa}

\begin{qa}{702}{Why could neither the controller nor the RL agent react to that loss?}
They could not see it. The drops happened in the fq\_codel leaf under each HTB class (CoDel drops at dequeue). The
collector read only the HTB class counters, which stayed at 0 in every one of those runs. The reward's loss term was
zero exactly when it mattered.
\ev{\Lthree}{a 40\,s disaster run with leaf counters reads $\approx$274 drops per tick where the class counters read 0.}
\end{qa}

\begin{qa}{703}{What else was structurally wrong?}
Three things.
\begin{enumerate}
\item \textbf{Direction.} The link is full duplex and HTB shapes each direction separately, but demand was summed over
both directions and compared with one direction's capacity. All healthcare applications are uplink (90\,Mb/s in a
disaster fits the 100\,Mb/s uplink), and education and general are downlink.
\item \textbf{Queue.} The default CoDel target (5\,ms) trimmed video bursts that healthcare's 50\,ms delay budget could
easily absorb.
\item \textbf{RL's lever.} RL only rescaled priority weights, which feed an allocator whose slices borrow idle capacity
anyway. Live, PPO chose multipliers of about 1.
\end{enumerate}
\end{qa}

\section{Module A: QoS visibility}

\begin{qa}{704}{What does the controller now measure per slice, and how?}
Every tick, per slice and interval:
\begin{itemize}
\item loss = queue drops / (sent + drops), over HTB class plus leaf qdisc, both directions;
\item queueing delay = standing queue / the slice's service rate;
\item delivered rate.
\end{itemize}
The result is \code{slice\_qos} in every snapshot (\code{telemetry/collector.py}). The SLA monitor, the demand
estimator, the reward, the RL state and the feedback log all read the same numbers.
\why{The bottleneck is the only contention point in the topology, so its own counters are a direct, cheap and complete
measurement of slicing-induced loss and delay.}
\choice{Kernel counters every tick.}
\alt{Per-flow measurements from live captures are more end-to-end, but cost a capture parse per tick. They remain the
offline ground truth (dataset KPIs).}
\trade{The queue estimate is a 1\,s snapshot, not a per-packet sojourn time.}
\ev{\Ltwo\ \Lthree}{\code{tests/test_telemetry_health.py}; module A regression.}
\end{qa}

\begin{qa}{705}{Does the reward now punish loss properly?}
Yes. A slice whose loss exceeds its SLA loss target (healthcare 1\,\%, education and agriculture 2\,\%, general 5\,\%)
now counts as an SLA violation ($-10$), in addition to the mean-loss term. The same function serves the simulators,
the live feedback log and the RL transitions.
\ev{\Ltwo}{\code{tests/test_reward.py::test_loss_above_sla_target_is_a_violation}.}
\end{qa}

\begin{qa}{706}{Did making loss visible break anything?}
Yes, and the regression caught it. The predictive backhaul detector treated \emph{any} drop as local congestion and
suppressed itself. With fq\_codel drops visible, steady TCP sawtooth drops muted it.
\begin{itemize}
\item First fix: use the core-bound slices' queueing delay (> 10\,ms) as congestion evidence, not drops. Proactive
switchover recovered from 1/3 to 3/4 runs.
\item Second fix: isolated single-tick TCP queue spikes still reset its 6-tick window, so congestion now suppresses
prediction only when it persists in at least half the window.
\end{itemize}
\ev{\Lthree}{switch-over before the cut in 5/5 repeated live runs; \code{test_isolated_queue_spikes_do_not_mask_a_backhaul_trend}.}
\end{qa}

\section{Module B: allocate for QoS targets, per direction}

\begin{qa}{707}{How does the QoS-aware allocator work?}
Per direction, with capacity $C$ each (\code{QoSAwareAllocator}):
\begin{enumerate}
\item floors (a 0.5\,Mb/s reverse floor in a slice's light direction);
\item \textbf{demands first}: if $\sum_s\max(f_s,d_s)\le C$ every slice gets its demand; otherwise slices are served in
descending context priority after reserving the others' floors, and those left short are listed as \emph{degraded};
\item \textbf{burst headroom} from the spare capacity, up to $d_s(h_s-1)$, first to slices currently failing their
SLA, then by priority;
\item leftover capacity to the slices that use this direction.
\end{enumerate}
\why{The goal is SLA compliance, not proportional fairness. When demand does not fit, which class degrades is a policy
decision, and it should be explicit and auditable.}
\alt{A weighted split (the audited \code{proposed}) degrades every class a little, implicitly.}
\trade{Strict priority degrades low-priority slices fully when the link is overcommitted.}
\ev{\Ltwo}{\code{tests/test_qos_aware.py}.}
\end{qa}

\begin{qa}{708}{What is the headroom loop?}
Each slice starts with headroom $h_s=1.1$. If its mean loss over the last interval exceeds its SLA loss target,
$h_s\leftarrow\min(1.6,\,1.15\,h_s)$ and the slice is marked at risk. After three intervals below 25\,\% of its target,
$h_s\leftarrow\max(1.1,\,0.97\,h_s)$. This is a per-slice AIMD-like loop closed on measured loss.
\ev{\Ltwo}{\code{test_headroom_rises_with_loss_and_relaxes_when_quiet}.}
\end{qa}

\begin{qa}{709}{Why derive per-slice CoDel parameters from the SLA?}
CoDel drops once the standing queue exceeds its target for an interval. Healthcare's one-way budget is 50\,ms, with
8\,ms of propagation, so a 12\,ms target with a 120\,ms interval trades a few ms of queueing for far fewer burst drops.
Education and agriculture get 25\,ms / 250\,ms; general keeps the 5\,ms / 100\,ms defaults.
\trade{Higher queueing delay and jitter inside the budget. The measured healthcare p95 stayed at 8--10\,ms (dev).}
\end{qa}

\begin{qa}{710}{Which demand does the direction-aware allocator plan for?}
The mean offered rate over the last control interval, per direction.
\begin{itemize}
\item Elastic probe-up (+25\,\%) applies only to directions where the service has TCP traffic, and only when a
standing queue proves saturation.
\item The first version followed per-second samples with a trend forecaster. On UDP video it estimated 228\,Mb/s for
a 95\,Mb/s flow (ramp-up trend plus probe-up on inelastic traffic), squeezed agriculture to its floor, and was replaced.
\end{itemize}
\ev{\Ltwo\ \Lthree}{\code{test_directional_estimate_probes_only_elastic_traffic_and_never_overshoots_peak}; development runs.}
\end{qa}

\begin{qa}{711}{What did module B achieve on the network during development?}
Healthcare SLA-met share:
\begin{itemize}
\item disaster: 0.79 $\rightarrow$ 0.93;
\item medical emergency: 0.67 $\rightarrow$ 0.87.
\end{itemize}
Healthcare loss fell from 0.7--1.4\,\% to $\approx$0.2\,\%, and weighted utility rose in both contexts.
Agriculture, which shares the uplink, fell from 1.0 to 0.94 and 0.75. Medical emergency offers 100\,Mb/s of VBR uplink
traffic on 100\,Mb/s, so not every class can be met there: the allocator now degrades agriculture explicitly (context
priority 15:1) instead of letting healthcare lose implicitly.
\ev{\Lthree}{\emph{(dev, $n{=}2$)}; confirmed or refuted in \S\ref{sec:qoseval}.}
\end{qa}

\section{Module C: react within a second, act ahead when possible}

\begin{qa}{712}{What is reactive, and what is predictive?}
\begin{itemize}
\item \textbf{Reactive.} A slice breaching its loss SLA, or with queueing delay above half its delay budget, triggers a
re-plan on the next 1\,s tick instead of waiting up to 5\,s. A 2-tick hold-down prevents oscillation. Explicitly
degraded slices do not retrigger, because re-planning cannot help them. Queue-based triggers ignore TCP slices: TCP
fills its own queue by design.
\item \textbf{Predictive.}
  \begin{itemize}
  \item a consistent demand ramp ($R^2\ge0.6$) is extrapolated two ticks ahead, capped at $1.5\times$ the window peak;
  \item a rising queue projected to cross the limit within two ticks triggers a re-plan;
  \item context changes predictable from the clock and calendar (the school bell) are pre-provisioned two ticks early.
  \end{itemize}
\end{itemize}
\honest{Alerts are not predictable. The scripted alert times in the timelines are never read: that would be cheating.}
\ev{\Ltwo\ \Lthree}{\code{tests/test_qos_adaptive.py};
\code{test_qos_adaptive_preprovisions_before_clock_driven_context_change}.}
\end{qa}

\begin{qa}{713}{Did module C measurably help in development?}
Not detectably. The mechanisms fire as designed (pre-provisioning two ticks before the bell, SLA-risk re-plans), but
on the two timelines \code{qos_adaptive} was within noise of \code{qos_aware} \emph{(dev, $n{=}2$)}. Burst headroom and
HTB borrowing already absorb most transients at a 5\,s cadence, and context changes are rare. The five-seed evaluation
decides (\S\ref{sec:qoseval}).
\end{qa}

\section{Module D: residual reinforcement learning}

\begin{qa}{714}{What does the redesigned RL control, and why that?}
A residual on the QoS-aware allocation: $a\in[-1,1]^4$ scales each slice's rate in its dominant direction by
$(1+0.3a_s)$. Each direction is renormalised to $C$, with floors kept.
\why{RL acts directly on the quantity that decides QoS (rates), and the heuristic is a strong, safe starting point.
$a{=}0$ \emph{is} the heuristic, so RL only has to learn corrections.}
\alt{Priority multipliers (the earlier design), whose effect HTB borrowing mostly cancels.}
\trade{RL replaces the hand-made headroom loop, so \code{rl\_residual} vs \code{qos\_adaptive} is a fair test of
``learned vs designed adaptation''.}
\ev{\Ltwo}{\code{tests/test_residual_rl.py} (floors and capacity hold for any action; $a{=}0$ reproduces the heuristic).}
\end{qa}

\begin{qa}{715}{What does it observe?}
State v2 (25 values), built by one function in both the simulator and the live controller:
\begin{itemize}
\item context one-hot;
\item per slice in its dominant direction: demand / $C$ and allocation / $C$;
\item loss / SLA loss target, and queueing delay / SLA delay, measured over the last interval;
\item the on-backup flag.
\end{itemize}
\end{qa}

\begin{qa}{716}{Why a new simulator, and how faithful is it?}
The old simulator had one pipe, no bursts, no AQM and no TCP behaviour, so it could not express the failure seen live.
The new one (\code{scheduler/env_dir.py}) models:
\begin{itemize}
\item both directions;
\item 50 sub-steps of 100\,ms per 5\,s step;
\item lognormal bursts;
\item HTB borrowing by priority;
\item CoDel burst tolerance $\approx$ rate $\times$ (target + interval);
\item TCP backing off instead of dropping.
\end{itemize}
Burstiness and HTB's borrowing efficiency were fitted to \NCalTicks{} live ticks with per-slice loss, minimising per-run mean
loss error: $\sigma_{\mathrm{hc}}{=}\NCalSigmaHc$, $\sigma_{\mathrm{agri}}{=}\NCalSigmaAg$, $\beta{=}\NCalBeta$.
\honest{The per-run RMSE is still \NCalRmseHc{} (healthcare) and \NCalRmseAg{} (agriculture) percentage points. The simulator guides training; only the network decides (Q725).}
\end{qa}

\begin{qa}{717}{What did PPO learn in simulation?}
Three seeds of 300k steps, best checkpoint chosen on validation seeds, evaluated on 30 paired held-out test episodes.
Reward against the heuristic ($a{=}0$): +\NResRelMin\,\% to +\NResRelMax\,\%, significant for \NResSigSeeds{} of the 3
seeds. SLA violations fell from
\NResViolHeuristic{} to \NResViolMin--\NResViolMax{} per episode. Small, and it did not carry over to the network (Q725).
\ev{\Ltwo}{\code{results/rl_training.json} (\code{residual}).}
\end{qa}

\begin{qa}{718}{Can it keep learning on the real network?}
\code{rl\_residual\_online} runs CEM over the residual action. It starts at $a{=}0$ (the heuristic), samples each live
decision and learns only from transitions whose allocation the kernel confirmed, with a checkpoint carried across runs.
\ev{\Lthree}{\code{test_residual_online_rl_closed_loop_on_live_network}: state v2 with measured QoS, exploring
actions, verified enforcement and policy updates from exactly the learnable transitions.}
\end{qa}

\section{Module E: SliceNet as the classifier for every flow}

\begin{qa}{719}{How does SliceNet classify on the go?}
In \code{slicenet\_all} mode there are no address rules: every UE flow starts best effort and is classified by SliceNet.
\begin{itemize}
\item Captures are read incrementally (only the new records each tick).
\item A flow is decided once it has been seen for 1\,s and 5 packets. Its features are expressed as a full 5\,s window
(rates over the observed span; counts scaled up).
\item Every 5\,s it is re-checked on a real full window. A different class predicted with confidence $\ge0.8$ on two
consecutive checks replaces the rule.
\item Rules are \code{tc flower} exact matches (hash lookup) with a handle per flow, so they can be replaced in place.
\item The confidence gate and best-effort fallback stay in force.
\end{itemize}
\ev{\Ltwo\ \Lthree}{\code{tests/test_online_slicenet.py}; \code{test_slicenet_classifies_every_ue_flow_without_address_rules}.
Development school run: 32 flows, median 1.57\,s to decision, 81\,\% correct at first decision and 87.5\,\% after two
(correct) reclassifications. Residual errors are education $\leftrightarrow$ general (both web-over-TCP).}
\end{qa}

\begin{qa}{720}{Why not classify inside a class (telemedicine vs ECG)?}
The evidence did not call for it. Inside the healthcare slice, fq\_codel already gives sparse flows (ECG, voice)
priority over bulky ones: their loss was 0\,\% while video lost 3--6\,\%. The fix was capacity and burst tolerance
for the slice, not a finer class.
\end{qa}

\section{Module F: evaluation in contended regimes}\label{sec:qoseval}

\begin{keybox}{Design of the evaluation}
\begin{itemize}
\item \textbf{Conditions.} Only conditions in which slices compete: disaster and medical emergency at 100\,Mb/s; school
and harvest on a 50\,Mb/s uplink; school and disaster at $1.5\times$ load; and the two timelines (a clock-driven school
bell and unpredictable alerts).
\item \textbf{Schedulers.} static\_equal, context\_unaware, proposed (the audited heuristic), qos\_aware (B),
qos\_adaptive (B+C) and rl\_residual (B+C+D), all on host classification.
\item \textbf{Size.} Five seeds per cell, \NQRuns{} runs, paired by (condition, seed) against \code{qos_adaptive}.
\item \textbf{Extra studies.} SliceNet everywhere (E) against address rules on the same controller; online residual RL
across consecutive runs.
\end{itemize}
The verdicts below come from the \textbf{confirmatory run v4 (seeds 15--19)}. Earlier rounds found defects; see Q732.
\end{keybox}

\begin{table}[!htbp]\centering\small
\begin{tabular}{@{}l ccc ccc ccc@{}}
\toprule
& \multicolumn{3}{c}{healthcare SLA share} & \multicolumn{3}{c}{agriculture SLA share} & \multicolumn{3}{c}{general SLA share}\\
\cmidrule(lr){2-4}\cmidrule(lr){5-7}\cmidrule(l){8-10}
condition & prop. & B+C & RL & prop. & B+C & RL & prop. & B+C & RL\\
\midrule
disaster          & \NQCondDisasterProposedHcSla & \NQCondDisasterQosAdaptiveHcSla & \NQCondDisasterResidualHcSla & \NQCondDisasterProposedAgSla & \NQCondDisasterQosAdaptiveAgSla & \NQCondDisasterResidualAgSla & \NQCondDisasterProposedGenSla & \NQCondDisasterQosAdaptiveGenSla & \NQCondDisasterResidualGenSla\\
medical emergency & \NQCondMedicalEmergencyProposedHcSla & \NQCondMedicalEmergencyQosAdaptiveHcSla & \NQCondMedicalEmergencyResidualHcSla & \NQCondMedicalEmergencyProposedAgSla & \NQCondMedicalEmergencyQosAdaptiveAgSla & \NQCondMedicalEmergencyResidualAgSla & \NQCondMedicalEmergencyProposedGenSla & \NQCondMedicalEmergencyQosAdaptiveGenSla & \NQCondMedicalEmergencyResidualGenSla\\
disaster $1.5\times$ & \NQCondDisasterOnepFivexProposedHcSla & \NQCondDisasterOnepFivexQosAdaptiveHcSla & \NQCondDisasterOnepFivexResidualHcSla & \NQCondDisasterOnepFivexProposedAgSla & \NQCondDisasterOnepFivexQosAdaptiveAgSla & \NQCondDisasterOnepFivexResidualAgSla & \NQCondDisasterOnepFivexProposedGenSla & \NQCondDisasterOnepFivexQosAdaptiveGenSla & \NQCondDisasterOnepFivexResidualGenSla\\
harvest 50\,Mb/s  & \NQCondHarvestFiveZeromProposedHcSla & \NQCondHarvestFiveZeromQosAdaptiveHcSla & \NQCondHarvestFiveZeromResidualHcSla & \NQCondHarvestFiveZeromProposedAgSla & \NQCondHarvestFiveZeromQosAdaptiveAgSla & \NQCondHarvestFiveZeromResidualAgSla & \NQCondHarvestFiveZeromProposedGenSla & \NQCondHarvestFiveZeromQosAdaptiveGenSla & \NQCondHarvestFiveZeromResidualGenSla\\
school 50\,Mb/s   & \NQCondSchoolFiveZeromProposedHcSla & \NQCondSchoolFiveZeromQosAdaptiveHcSla & \NQCondSchoolFiveZeromResidualHcSla & \NQCondSchoolFiveZeromProposedAgSla & \NQCondSchoolFiveZeromQosAdaptiveAgSla & \NQCondSchoolFiveZeromResidualAgSla & \NQCondSchoolFiveZeromProposedGenSla & \NQCondSchoolFiveZeromQosAdaptiveGenSla & \NQCondSchoolFiveZeromResidualGenSla\\
school $1.5\times$ & \NQCondSchoolOnepFivexProposedHcSla & \NQCondSchoolOnepFivexQosAdaptiveHcSla & \NQCondSchoolOnepFivexResidualHcSla & \NQCondSchoolOnepFivexProposedAgSla & \NQCondSchoolOnepFivexQosAdaptiveAgSla & \NQCondSchoolOnepFivexResidualAgSla & \NQCondSchoolOnepFivexProposedGenSla & \NQCondSchoolOnepFivexQosAdaptiveGenSla & \NQCondSchoolOnepFivexResidualGenSla\\
rural day (timeline) & \NQCondRuralDayProposedHcSla & \NQCondRuralDayQosAdaptiveHcSla & \NQCondRuralDayResidualHcSla & \NQCondRuralDayProposedAgSla & \NQCondRuralDayQosAdaptiveAgSla & \NQCondRuralDayResidualAgSla & \NQCondRuralDayProposedGenSla & \NQCondRuralDayQosAdaptiveGenSla & \NQCondRuralDayResidualGenSla\\
medical night (timeline) & \NQCondMedicalNightProposedHcSla & \NQCondMedicalNightQosAdaptiveHcSla & \NQCondMedicalNightResidualHcSla & \NQCondMedicalNightProposedAgSla & \NQCondMedicalNightQosAdaptiveAgSla & \NQCondMedicalNightResidualAgSla & \NQCondMedicalNightProposedGenSla & \NQCondMedicalNightQosAdaptiveGenSla & \NQCondMedicalNightResidualGenSla\\
\bottomrule
\end{tabular}
\caption{Share of flow windows meeting their application's SLA, per class (mean over 5 seeds; confirmatory run v4).
prop.\ = \capcode{proposed} (audited heuristic), B+C = \capcode{qos_adaptive}, RL = \capcode{rl_residual}. Education met its SLA
in nearly every window under every scheduler and is omitted. Generated from \capcode{results/tables/qos_summary.csv}.}
\label{tab:qos}
\end{table}

\FigQosHealthcare

\begin{qa}{721}{Did the QoS-targeted control make the critical class meet its SLA where it previously did not?}
Largely yes. Over the 40 paired runs, healthcare's SLA-met share rises by \NQDiffProposedHcSla{} against
\code{proposed} ($p=\NQPProposedHcSla$, significant after Bonferroni). The healthcare-loss difference (QoS-targeted $-$ proposed) is
\NQDiffProposedHcLoss{} percentage points.
\begin{itemize}
\item Disaster: \NQCondDisasterProposedHcSla{} $\rightarrow$ \NQCondDisasterQosAdaptiveHcSla.
\item Medical emergency: \NQCondMedicalEmergencyProposedHcSla{} $\rightarrow$ \NQCondMedicalEmergencyQosAdaptiveHcSla.
\item Night-time medical alert: \NQCondMedicalNightProposedHcSla{} $\rightarrow$ \NQCondMedicalNightQosAdaptiveHcSla.
\end{itemize}
It is not a full fix. At $1.5\times$ disaster load (135\,Mb/s of healthcare on a 100\,Mb/s uplink) no allocation can
meet everyone, and healthcare stays near \NQCondDisasterOnepFivexQosAdaptiveHcSla.
\ev{\Lthree}{\code{results/tables/qos_stats.csv}.}
\end{qa}

\begin{qa}{722}{Then why is H11 ``not supported''?}
H11 was stated on the \emph{unweighted mean} over the four classes, and on that metric \code{qos_adaptive} is
\NQDiffProposedMeanSla{} below \code{proposed} ($p=\NQPProposedMeanSla$). The gain for healthcare is paid for by
agriculture, which shares the uplink and has context priority 1 against healthcare's 10--15 in the emergencies:
\begin{itemize}
\item disaster: \NQCondDisasterProposedAgSla{} $\rightarrow$ \NQCondDisasterQosAdaptiveAgSla;
\item medical emergency: \NQCondMedicalEmergencyProposedAgSla{} $\rightarrow$ \NQCondMedicalEmergencyQosAdaptiveAgSla.
\end{itemize}
General also gives a little up at 50\,Mb/s. This is the declared priority order working as specified, not a defect;
it was deliberately \emph{not} tuned away after it was seen. The worst-off class is \NQDiffProposedMinSla{} better
($p=\NQPProposedMinSla$, not significant), weighted utility is unchanged ($p=\NQPProposedUtil$), and against static
slicing the worst-off class is \NQDiffStaticMinSla{} better ($p=\NQPStaticMinSla$).
\honest{Whether that trade is right is a policy question for the operator. The framework makes it explicit and
measurable; it does not make the trade disappear.}
\end{qa}

\begin{qa}{723}{Where does the QoS-targeted control clearly win, and where does it not help?}
\begin{itemize}
\item \textbf{Clear win: overloaded uplink with a protected critical class.} Harvest on 50\,Mb/s: worst-off class
\NQCondHarvestFiveZeromProposedMinSla{} $\rightarrow$ \NQCondHarvestFiveZeromQosAdaptiveMinSla, and agriculture
\NQCondHarvestFiveZeromProposedAgSla{} $\rightarrow$ \NQCondHarvestFiveZeromQosAdaptiveAgSla. Direction-aware
allocation stops wasting uplink guarantees on downlink services.
\item \textbf{No help: school.} School is overloaded on the \emph{downlink} by TCP traffic: both education and general
are elastic, and the weighted heuristic already does about as well.
\item \textbf{Trade-off: emergencies.} As in Q722.
\end{itemize}
\end{qa}

\begin{qa}{724}{Did reacting within a second and predicting help (module C, H13)?}
No measurable effect: \code{qos_adaptive} vs \code{qos_aware} differs by \NQDiffQosAwareMeanSla{} in mean class SLA
($p=\NQPQosAwareMeanSla$). The triggers fired \NQReactions{} times across the \NQRuns{} runs and did what they were
designed to do. At this load and timescale, burst headroom plus HTB borrowing already absorb the transients that a 1\,s
reaction could catch, and the predictable context changes are few.
\end{qa}

\begin{qa}{725}{Did residual RL improve on the heuristic on the network (H12)?}
No. \code{rl_residual} is within \NQDiffResidualMeanSla{} of \code{qos_adaptive} on mean class SLA
($p=\NQPResidualMeanSla$) and \NQDiffResidualMinSla{} on the worst class ($p=\NQPResidualMinSla$). In simulation it
gains +\NResRelMin{} to +\NResRelMax\,\% reward (significant for \NResSigSeeds{} of 3 seeds) and cuts SLA violations
from \NResViolHeuristic{} to \NResViolMin--\NResViolMax{} per episode. That gain does not survive the
simulation-to-network gap at the size we can measure with 40 paired runs.
\honest{A learned controller that only matches a well-designed heuristic is a negative result for RL here, not a
contribution. It is safe, though: it never breaks a floor or the anti-starvation guard (Q732), and \NQEnforcementFailed{}
of its allocations failed kernel verification.}
\end{qa}

\begin{qa}{726}{Did the online residual learner work on the real loop?}
Mechanically yes. Over \NQOnlineRuns{} consecutive live runs it made \NQOnlineTransitions{} verified transitions
(\NQOnlineTransMin--\NQOnlineTransMax{} per run, \NQOnlineFailed{} enforcement failures) with a checkpoint carried
across runs, and mean class SLA was \NQOnlineMeanSla.
\honest{About 25 transitions per run is far too few to show learning. It demonstrates a coupled, safe online loop, not
convergence.}
\end{qa}

\begin{qa}{727}{Can SliceNet replace address rules for every flow (H14)?}
Yes, within the pre-stated margin. With every UE flow classified by SliceNet (no address rules), mean class SLA is
\NQSnAllMeanSla{} against \NQSnHostMeanSla{} with address rules ($p=\NQSnPMeanSla$; margin 0.02). Decisions take a median
of \NQSnDecisionS\,s, with \NQSnAccuracy{} accuracy at first decision and \NQSnFinalAccuracy{} after confirmed
reclassification. The residual errors are education $\leftrightarrow$ general, whose web-over-TCP traffic looks alike.
\end{qa}

\begin{qa}{728}{Are the control actions actually enforced?}
Yes. Across the \NQRuns{} contended-regime runs, \NQChangesApplied{} allocation changes were applied per direction and
read back from the kernel; \NQEnforcementFailed{} failed verification.
\end{qa}

\section{How the evaluation was kept honest}

\begin{qa}{732}{The QoS schedulers were evaluated four times. Isn't that tuning on the test set?}
Each round used fresh seeds, and only \emph{defects} were fixed: behaviour that violated a stated design requirement.
Each fix has a regression test. All rounds are reported (\code{results/evidence/qos_eval_v1}--\code{v3}).
\begin{description}
\item[v1 (seeds 0--4)] Strict priority on bare demand starved best-effort TCP (general SLA 0.0 at 50\,Mb/s), and the
overload branch ignored burst headroom. Fixed: an anti-starvation guard (20\,\% of demand) and headroom applied under
overload.
\item[v2 (seeds 5--9)] A starved TCP slice sends less, so its measured demand, and with it the guard, collapsed
(general SLA 0.05). Fixed: the guard of an elastic slice is anchored on the context's design demand.
\item[v3 (seeds 10--14)] The residual RL could take capacity from a slice down to its plain floor, overriding the
guard (general SLA 0.04 under \code{rl_residual}). Fixed: the allocator's protected minimum bounds the residual, live
and in the simulator.
\item[v4 (seeds 15--19)] The confirmatory run reported here.
\end{description}
What was \emph{not} changed is equally important. Agriculture's degradation in the emergencies follows the declared
context priority; it was reported, not tuned away.
\end{qa}

\begin{qa}{733}{So, can SliceNet and RL slice the network so that the four classes meet their QoS on the go?}
Partly, and the parts are now measured.
\begin{itemize}
\item \textbf{SliceNet} can identify every flow on the fly (about 1.5\,s) and drive slicing without address rules, at no
measurable QoS cost (H14).
\item \textbf{The QoS-targeted control}, which now sees loss, allocates per direction, keeps burst headroom and never
starves a class, delivers the critical class its SLA far more often in the emergencies and relieves the worst-off class
when an uplink is overloaded. Where the link cannot carry everything, it degrades in the declared priority order rather
than meeting all four classes.
\item \textbf{RL}, redesigned to act on rates with measured QoS in its state, is safe and live-coupled, but does not
beat the heuristic it corrects.
\item \textbf{Prediction} (ramps, queues, clock-driven context) adds nothing measurable here.
\end{itemize}
The honest summary for the panel: the gains came from \emph{measuring the right thing and allocating for it}, not from
learning.
\end{qa}

